\section{Comparing two cut-offs and identifying the limit}\label{app:closing}
The preceding estimates hold for each fixed admissible cut-off. Their thresholds for \(n\) may depend
on the cut-off, so passage to the endpoint still requires justification. The quantity they approximate
on the initial board is the same for all cut-offs. Let \(a_n=nR_Q(n)\), \(a_*=F(1)-1/12\),
\(H=P_T+J-S\), and let \(b>0\) be an admissible cut-off. There are admissible cut-offs arbitrarily
close to zero. At a positive cut-off define \(I(b)=\int_b^1 k(x)x^{-2}\,dx\) and
\(\Phi(b)=a_*+I(b)+H\). The three boundary identifications give \(\eta(b)=\limsup_n|B_n(b)-H|\to0\).
The interior identity holds on \([b,1]\). Combining it with \eqref{eq:signed-mean}, the boundary
expansion with error \(O(n/j^2)\) and the limit \(n g_T(M_n)\to a_*\), we obtain
\[
 \limsup_n|a_n-\Phi(b)|\le C_\delta b+\eta(b)=:\varepsilon(b).
\]
\begin{lema}[Comparison of two cut-offs]\label{lem:dos-cortes}
Let \((a_n)\) be a real sequence and, for each admissible cut-off \(b\), let \(\Phi(b)\in\mathbb R\)
satisfy \(\limsup_n|a_n-\Phi(b)|\le\varepsilon(b)\), where \(\varepsilon(b)\to0\) as the cut-off
decreases to zero. Then \(\Phi(b)\) has a finite limit \(\Phi_0\) along the admissible cut-offs and
\(a_n\to\Phi_0\).
\end{lema}
\begin{proof}
For two fixed admissible cut-offs we compare \(\Phi(b),\Phi(b')\) with the same \(a_n\), taking \(n\)
sufficiently large for both estimates. The triangle inequality and the limit superior give
\(|\Phi(b)-\Phi(b')|\le\varepsilon(b)+\varepsilon(b')\). The Cauchy criterion holds along the
admissible cut-offs and \(\Phi(b)\to\Phi_0\in\mathbb R\). For each cut-off,
\(\limsup_n|a_n-\Phi_0|\le\varepsilon(b)+|\Phi(b)-\Phi_0|\). Since the left-hand side is independent
of \(b\), letting the cut-off decrease to zero gives \(a_n\to\Phi_0\).
\end{proof}
In particular, \(I(b)\) has a finite limit along these cut-offs. Throughout this argument we have used
the physical profile, the stopped laws and their estimates at a fixed positive cut-off.

\begin{proof}[Conclusion of the proof of Theorem~\ref{teo:principal} and Corollary~\ref{cor:derivada}]
To identify the limit with the ordinary improper integral, it remains to decide whether \(k(x)/x^2\)
can have a term proportional to \(1/x\). The zeroth-order profile identity gives \(k(0)=0\). The
projected line resolvent is analytic near zero because the diagonal Schur block has the form
\(\Lambda_u+x\mathcal N\), the vacancy has a positive minimum and the axial block has a spectral gap.
The inverse and its kernel sections remain bounded on interior complex disks. Dominated integration in
the finite contractions of the densities thus gives analyticity of \(k,G\). This regularity holds for
any value of \(k'(0)\). On a sufficiently small interval \([0,x_0]\), write \(k(x)=k'(0)x+x^2 h(x)\),
with \(h\) analytic and bounded. For \(0<b<x_0\),
\[
 I(b)=I(x_0)+k'(0)\log(x_0/b)+\int_b^{x_0}h(x)\,dx .
\]
The last integral remains bounded. Convergence of \(I(b)\) along the admissible cut-offs forces
\(k'(0)=0\). From \(k'=d'-xG''\) we deduce \(d'(0)=0\) for the physical profile. The singularity of
the integrand is then removable and the improper integral converges absolutely. Consequently, the
value of \(\Phi_0\) is
\[
 \lim_{n\to\infty}n\{\log Q(n)-\log Q_{\rm pred}(n)\}
   =F(1)-\tfrac1{12}+\int_0^1\frac{k(x)}{x^2}\,dx+P_T+J-S .
\]
\end{proof}
